\subsection{Monte Carlo simulation}

\begin{qfmethod}[Monte Carlo simulation: sampling from a specified model, the estimation target, seeds, standard error, the square-root rule, control variates and discretisation bias]\label{meth:monte-carlo-simulation}
\qffield{Statement} A Monte Carlo valuation samples paths of the variables from a specified model, computes the discounted payoff on each, and averages. With $\bar f$ the mean and $\omega$ the standard deviation of the discounted payoffs over $M$ trials, the standard error is $\omega/\sqrt M$ and a 95 percent confidence interval for the value is $\bar f\pm1.96\,\omega/\sqrt M$; uncertainty falls with the square root of the number of trials, so doubling the accuracy needs four times the trials \citep[chap.~21, sec.~21.6, pp.~492--494]{s7-hulloptions}. In Hull's example, 1{,}000 trials give an option value of 4.98 with a standard error of 0.24, and the closed-form price, 4.817, lies inside the interval (Example~21.8, p.~494).

\emph{Control variates.} When a similar derivative B has an analytic value $f_B$, both are simulated on the same random numbers and the estimate of A is corrected to $f_A=f_A^*-f_B^*+f_B$ (sec.~21.7, eq.~(21.20), p.~496). \emph{Discretisation against sampling error.} Simulating $\ln S$ with constant parameters is exact at any step, while an Euler step of $dS$ is exact only as $\Delta t\to0$ (p.~491 and n.~13). In general the time step contributes an error of order $\delta t$ and the finite number of paths one of order $N^{-1/2}$. For a fixed budget proportional to $N/\delta t$, Wilmott balances them with $N$ of order $K^{2/3}$ and $\delta t$ of order $K^{-1/3}$ \citep[chap.~80, secs.~80.4--80.6, pp.~1266--1268]{wide-wilmott-pwoqf}.

\qffield{Notation} Target: $\theta=E[g(Y)]$ under the specified distribution of $Y$, for example a discounted payoff in a risk-neutral model or a probability $P(Y\in A)=E[\mathbf 1_A(Y)]$, whose standard error is $\sqrt{\hat p(1-\hat p)/M}$. Seed: the starting state of the pseudo-random generator, which makes a run reproducible. Interpretation: with a control $X$ of known mean, $Y-b(X-E X)$ has the same mean as $Y$ and the smallest variance at $b^*=\mathrm{Cov}(Y,X)/\mathrm{Var}(X)$, where its variance is $(1-\rho^2)\mathrm{Var}(Y)$. Hull's correction is the case $b=1$. Estimating $b^*$ from the same draws adds a bias of order $1/M$, negligible at these sizes.

\qffield{Conditions} The estimate is only as good as the model: Monte Carlo gives the expectation under the assumed distribution, not under the true one. The standard-error formula assumes independent trials with finite payoff variance; dependent draws, such as overlapping paths, need other error estimates. Variance-reduction methods such as antithetic variables change the precision per trial, not the target (Hull, sec.~21.7). Seeds should be fixed before the run and reported, not chosen after seeing results. The standard error measures sampling error only: discretisation bias does not shrink with more paths, so it is measured separately, for example against the exact solution on common Brownian paths.

\qffield{Practice} Example, original. A stock at 100 follows geometric Brownian motion with volatility 20 percent; the risk-free rate is 3 percent; a one-year call struck at 105 has a closed-form value of 7.1281. With 100{,}000 draws the estimate is 7.1161 and the standard error 0.0396, so the closed form lies 0.3 standard errors away. The probability that the price ends below 90 is 0.2821 by simulation against 0.2820 exactly (standard error 0.0014). A quarter of the draws gives about twice the standard error, and antithetic pairs reduce it further.

With the discounted terminal price as control (known mean 100), the estimated coefficient is 0.542 and the correlation 0.875. The variance falls by $1/(1-\rho^2)=4.25$, and the estimate becomes 7.1289 with standard error 0.0192, against 0.0396 without the control. \emph{Discretisation}: with Euler steps of $dS$ on the same 400{,}000 Brownian paths as the exact solution, the value is biased by $-0.125$, $-0.054$ and $-0.025$ at 2, 4 and 8 steps. These biases are many times their standard errors (0.002 or less) and halve with the step, whereas more paths would leave them unchanged.

\qffield{Pitfalls} Reporting a Monte Carlo figure without its standard error, number of trials and seed. Reading agreement within sampling error as validation of the model. Comparing two simulated strategies on different random draws when common random numbers would sharpen the comparison.

\qffield{See also} \hyperref[def:financial-distributions]{distributions} (lognormal prices); \hyperref[mod:black-scholes-merton]{Black--Scholes--Merton}; \hyperref[meth:var-es-estimation-methods]{VaR and ES estimation} (Monte Carlo for risk); \hyperref[res:ito-lemma-geometric-brownian-motion]{It\^o's lemma} (exact and Euler sampling); \hyperref[meth:least-squares-monte-carlo]{least-squares Monte Carlo} (early exercise); \hyperref[meth:asian-option-valuation]{Asian options} (a geometric-average control); \hyperref[meth:monte-carlo-greeks]{Monte Carlo Greeks} (sensitivities from the same simulation).

\end{qfmethod}
